\section{Frobenius quotients and canonical normalization}
\label{rev:truncations}\label{sec:general-pullbacks}

A first jet governs transport of a complex identity. Arithmetic
transport depends on a different object: the quotient of a normalized
series by its Frobenius pullback. The quotient is functorial; its raw
Taylor truncation need not be. This distinction both explains the
finite transformations and gives a precise rigidity theorem for the
standard normalization.

\subsection{Lucas congruences as a polynomial Frobenius structure}
\label{subsec:C-canonical-lucas}
For $h\in1+t\mathbb F_p[[t]]$, write
\[
 \mathcal Q_p(h)=\frac{h(t)}{h(t^p)},\qquad
 [h]_{<p}=\sum_{n<p}[t^n]h\,t^n.
\]
The coefficient congruences $h_{pa+b}=h_ah_b$ are equivalent to
\begin{equation}\label{rev:eq:lucas-quotient}
 \mathcal Q_p(h)=[h]_{<p}.
\end{equation}
Indeed $h(t^p)=1+O(t^p)$, so this equality says exactly that
$\mathcal Q_p(h)$ is a polynomial of degree less than $p$.
For normalized $g$ and $r\in t\mathbb F_p[[t]]$, Frobenius gives
the commutative diagram
\[
\begin{tikzcd}[column sep=large]
 h \arrow[r,mapsto,"h\mapsto g(h\circ r)"]
   \arrow[d,mapsto,"\mathcal Q_p"']
 &g(h\circ r)\arrow[d,mapsto,"\mathcal Q_p"]\\
 \mathcal Q_p(h)\arrow[r,mapsto,"Q\mapsto g^{1-p}(Q\circ r)"']
 &g^{1-p}(\mathcal Q_p(h)\circ r).
\end{tikzcd}
\]
Thus the obstruction to transporting Lucas congruences is simply
failure of the lower-right expression to remain a polynomial of
controlled degree.

\begin{proposition}[The standard Frobenius quotient]
\label{prop:C-standard-lucas}
Put $u_\ell(n)=R_\ell^n c_{s_\ell}(n)$, as in \eqref{eq:series}, and
$k_\ell=6,4,3,2$ for $\ell=1,2,3,4$. For every prime $p$,
\[
 u_\ell(pa+b)\equiv u_\ell(a)u_\ell(b)\pmod p
 \qquad(a\ge0,\ 0\le b<p).
\]
For $p>k_\ell$, the polynomial
$T_{\ell,p}(t)=\sum_{n<p}u_\ell(n)t^n$ has degree
$d_\ell(p)=\lfloor(p-1)/k_\ell\rfloor$. Consequently
\begin{equation}\label{eq:C-standard-lucas-functional}
 G_\ell(t)=T_{\ell,p}(t)G_\ell(t^p),\qquad
 G_\ell(t)=\sum_{n\ge0}u_\ell(n)t^n,
\end{equation}
in $\mathbb F_p[[t]]$. The same assertions hold for
$F_{s_\ell}(z)=G_\ell(z/R_\ell)$ when $p\nmid R_\ell$.
\end{proposition}
\begin{proof}
The factorial ratios have the integral presentations
\[
 \binom{6n}{3n}\binom{3n}{n,n,n},\quad
 \binom{4n}{n,n,n,n},\quad
 \binom{3n}{n,n,n}\binom{2n}{n},\quad
 \binom{2n}{n}^{\!3}.
\]
If $b<p/k_\ell$, their lowest base-$p$ digits have no carries, so
Lucas's multinomial theorem separates the digit $b$ from the higher
digits. If $b\ge p/k_\ell$, Legendre's valuation formula makes both
sides zero: its periodic summands on $[0,1)$ are, respectively,
\[
 \lfloor6x\rfloor-\lfloor3x\rfloor,\quad \lfloor4x\rfloor,
 \quad\lfloor3x\rfloor+\lfloor2x\rfloor,\quad3\lfloor2x\rfloor,
\]
all nonnegative and positive on $[1/k_\ell,1)$. This proves Lucas.
For $0\le n<p/k_\ell$ all factorial arguments are below $p$,
so the coefficient is a unit; the remaining coefficients below $p$
vanish. This proves the degree formula. Digit decomposition proves
\eqref{eq:C-standard-lucas-functional}, and Fermat's theorem permits
the rescaling by $R_\ell$.
\end{proof}

\begin{proposition}[Rational gauge rigidity]
\label{prop:C-lucas-rational-rigidity}
Suppose $f\in1+t\mathbb F_p[[t]]$ has polynomial Frobenius quotient
$T$ with $1\le\deg T<p-1$. If $g\in\mathbb F_p(t)$ is regular at
zero, $g(0)=1$, and $gf$ is Lucas, then $g=1$.
Hence a nonconstant rational gauge over $\Q$, normalized at zero,
destroys the Lucas property of each standard series at every
sufficiently large prime.
\end{proposition}
\begin{proof}
Write $g=P/Q$ in lowest terms. Polynomiality of
\begin{equation}\label{eq:C-lucas-gauge-covariance}
 \mathcal Q_p(gf)=Tg^{1-p}=\frac{TQ^{p-1}}{P^{p-1}}
\end{equation}
forces $P^{p-1}\mid T$, hence $P$ is constant. The degree bound
then forces $\deg Q=0$, and normalization gives $g=1$.
A fixed nonconstant rational function remains nonconstant and in
lowest terms outside finitely many primes, which proves the last
assertion.
\end{proof}

The covariance identity also distinguishes the quotient from its
truncation. If $g(t_0)=1$ and $g'(t_0)=0$ at a regular
$t_0\in\mathbb F_p$, then the first jets of $\mathcal Q_p(gf)$ and
$T$ agree there. Their raw truncations can disagree: truncation
discards the rational Frobenius tail. This is an identity in
characteristic $p$, not a comparison with a complex period value.

\subsection{Finite transport}
Put $\delta=t\partial_t$, and retain the branch and nondegeneracy
conditions of \eqref{eq:general-pullback-transport}.

\begin{lemma}[Polynomial transport criterion]
\label{lem:finite-frobenius-pullback}
Let $p>3$, $g\in1+t\Z_p[[t]]$, $r\in t\Z_p[[t]]$, and
$H=gF_s(r)$. If the reduction of
$\Lambda_p=g^{1-p}T_p(r)$ is a polynomial of degree less than $p$,
then
\begin{equation}\label{eq:finite-frobenius-pullback}
 [H]_{<p}\equiv\Lambda_p\pmod p.
\end{equation}
Writing $u=\delta\log g$ and $v=\delta\log r$, one also has
\begin{equation}\label{eq:finite-weighted-pullback}
 (A+B\delta)[H]_{<p}
 \equiv g^{1-p}\bigl((A+Bu)T_p(r)+Bv\,rT'_p(r)\bigr)\pmod p.
\end{equation}
The weighted equality is evaluated only where its expressions are
regular.
\end{lemma}
\begin{proof}
The standard Lucas identity and the diagram identify
$\Lambda_p$ with $\mathcal Q_p(H)$. The degree condition and
\eqref{rev:eq:lucas-quotient} prove the first equality. The second
is its derivative, since $\delta\log(g^{1-p})=u$ in characteristic
$p$.
\end{proof}

\begin{proposition}[Three finite transformations]
\label{prop:three-finite-transformations}
The quadratic, Domb, and level-nine transformations in
Appendix~\ref{rev:sec:elliptic-extensions} have the following exact
reductions for every $p>3$:
\begin{align}
 [\mathcal G]_{<p}(t)
 &\equiv(1+Kt)^{p-1}T_p\!\left(\frac{Rt}{(1+Kt)^2}\right),
 &&R=4K,\label{eq:finite-quadratic}\\
 [H_{\mathrm{Domb}}]_{<p}(t)
 &\equiv(1-4t)^{p-1}T_{p,1/3}\!\left(\frac{108t^2}{(1-4t)^3}\right),
 \label{eq:finite-domb}\\
 [H_9]_{<p}(t)
 &\equiv D(t)^{(p-1)/2}T_{p,1/6}\!\left(\frac{1728t^3}{D(t)^3}\right),
 &&D(t)=1-6t-27t^2.\label{eq:finite-nine}
\end{align}
Here $H_{\mathrm{Domb}}$ and $H_9$ denote the respective generating
functions in that appendix, and $T_{p,s}=[F_s]_{<p}$.
All right sides are polynomials of degree less than $p$, and their
weighted forms use the same logarithmic derivatives as the analytic
transport.
\end{proposition}
\begin{proof}
The degree of $T_{p,s}$ is $\lfloor p/k_\ell\rfloor$. Thus the
$n$th terms on the right have the forms
\[
 t^n(1+Kt)^{p-1-2n},\quad
 t^{2n}(1-4t)^{p-1-3n},\quad
 t^{3n}D^{(p-1)/2-3n},
\]
respectively. Every exponent is nonnegative, and the corresponding
degrees are at most $p-1$. The lemma applies to the three defining
germ identities.
\end{proof}

\begin{theorem}[Arithmetic classification after finite transport]
\label{thm:additional-raw-density}
For the ordinarily convergent quadratic-companion, Domb, and
level-nine series, fix an integer denominator $C\ne0$ and primitive
integer weights $(A,B)$. Omitting finitely many bad primes, put
\[
 Z=\{p:(A+B\delta)[H]_{<p}(1/C)=0\pmod p\}.
\]
The following conditions are equivalent:
\begin{enumerate}
\item $Z$ has lower natural prime density greater than $3/4$;
\item $Z$ contains every sufficiently large prime;
\item the data, up to simultaneous sign, are one of the five
quadratic-companion, two Domb, or two level-nine CM rows.
\end{enumerate}
\end{theorem}
\begin{proof}
At a regular convergent parameter the coefficient map is
\[
                 (A,B)\longmapsto(A+Bu,Bv),
\]
with rational entries and nonzero determinant. The proposition
identifies its standard weighted truncation with the original one,
up to a nonzero factor. Their zero sets therefore differ by a finite
set. Theorem~\ref{thm:weighted-density-cm} forces CM, the exact CM
classification determines the parameter, and
Proposition~\ref{prop:cm-weighted-slope} selects the complementary CM
eigenline. The period characterization of that line and invertible
analytic transport identify precisely the listed weights. Those
weights give cofinite zeros, proving both directions.

The only convergent parameters omitted by this regular argument are
$C=K$ for quadratic companions and $C=16$ for Domb. Convergence
forces $B=0$; finite transport then gives the zero condition
$T_p(1)=U_p(1/2)^2=0$. The elliptic fiber at $x=1/2$ is smooth and
CM, so Deuring's theorem gives density $1/2$. These endpoints
contribute no further row.
\end{proof}
This is an unconditional arithmetic theorem. It does not infer its
prime-density hypothesis from a complex identity at an unknown
non-CM fiber.

\subsection{Algebraic gauges and the half-density threshold}
\label{subsec:C-majority-lucas-gauges}
Rational rigidity admits one algebraic exception: an inverse square
root of a linear polynomial. The degree of the function-field
extension explains why a majority of primes already detects it.

\begin{theorem}[Classification of algebraic Lucas gauges]
\label{thm:C-majority-lucas-gauges}
Let $s\in\{1/6,1/4,1/3,1/2\}$ and let $g\in\Q[[z]]$ be algebraic
over $\Q(z)$ with $g(0)=1$. Let $\mathcal P_g$ be the primes at
which $gF_s$ is integral and Lucas. The following are equivalent:
\begin{enumerate}
\item $\mathcal P_g$ has upper natural prime density greater than $1/2$;
\item $\mathcal P_g$ contains every sufficiently large prime;
\item $g=(1-az)^{-1/2}$ for some $a\in\Q$, on the branch at zero.
\end{enumerate}
Under these conditions, either $g'(0)=0$, or $g(z_0)=1$ at a nonzero
regular point of an analytic continuation, forces $g=1$.
\end{theorem}
\begin{proof}
Let $d=[\Q(z,g):\Q(z)]$. The inclusion
$\Q(z,g)\subset\Q((z))$ makes this function field regular over
$\Q$: a Laurent series algebraic over $\Q$ has derivative zero
and is a rational constant. A primitive polynomial for $g$ is
therefore absolutely irreducible. Outside finitely many primes its
reduction remains absolutely irreducible with the same degrees.
For completeness, reducible forms of fixed degree form a closed
set: they are the finite union of the proper images of multiplication
maps between projective spaces of forms. A polynomial avoiding that
set in characteristic zero meets it at only finitely many primes.
Eisenstein's algebraic-power-series theorem also makes $g$
integral outside finitely many primes. Hence
\begin{equation}\label{eq:C-gauge-degree-reduction}
 [\mathbb F_p(z,g_p):\mathbb F_p(z)]=d
\end{equation}
for every sufficiently large $p$.

At a Lucas prime, write $T_p=\mathcal Q_p(F_s)$ and
$U_p=\mathcal Q_p(gF_s)$. Covariance gives
\begin{equation}\label{eq:C-algebraic-gauge-kummer}
                    g_p^{p-1}=T_p/U_p.
\end{equation}
Since $\mu_{p-1}\subset\mathbb F_p$, the extension generated by
$g_p$ is cyclic of degree dividing $p-1$: its automorphisms act by
multiplying $g_p$ by $(p-1)$st roots of unity. Thus
$\mathcal P_g$ is, up to finitely many exceptions, contained in
$p\equiv1\pmod d$. The prime number theorem in arithmetic
progressions gives this progression density $1/\varphi(d)$.
The strict half-density hypothesis forces $d\le2$.

If $d=1$, rational gauge rigidity applies. If $d=2$, the nontrivial
automorphism sends $g_p$ to $-g_p$ at every sufficiently large
Lucas prime. The rational trace of $g$ therefore vanishes modulo
infinitely many primes and is zero. Consequently $g^2=R\in\Q(z)$.
It remains to bound the divisor of $R$.

Every root of $T_p$ has multiplicity at most two. Indeed its degree
$d_p=\lfloor(p-1)/k_\ell\rfloor$ is at most $(p-1)/2$, and
\[
 L_sT_p=0,\qquad
 L_s=\thetaop^3-z(\thetaop+1/2)(\thetaop+s)(\thetaop+1-s),
\]
because derivatives kill the Frobenius factor of $F_s$. At an
ordinary point, a root of multiplicity $r\ge3$ contradicts the
nonvanishing indicial coefficient $r(r-1)(r-2)$, since $r<p$.
There is no root at zero. At $z=1$ the indicial expression is
$r(r-1)(r-1/2)$; its additional positive residue $(p+1)/2$ exceeds
$d_p$, so the multiplicity there is at most one.

Write $R=P/Q$ in lowest terms, with $P(0)=Q(0)\ne0$. The identity
\[
                     U_p=T_p(Q/P)^{(p-1)/2}
\]
forces $P$ to be constant: otherwise, for large $p$, a zero of $P$
would require a zero of $T_p$ of multiplicity at least $(p-1)/2$.
Normalize $R=1/Q_0$, $Q_0(0)=1$. The remaining divisor is bounded
at infinity:
\[
 \deg U_p=d_p+\frac{p-1}{2}\deg Q_0\le p-1.
\]
Since $d_p>0$, this forces $\deg Q_0\le1$, proving (3).

Conversely, $g=(1-az)^{-1/2}$ gives
\[
 \mathcal Q_p(gF_s)=T_p(z)(1-az)^{(p-1)/2},
\]
a polynomial of degree at most $p-1$ at every sufficiently large
prime. This proves (2), hence (1). Finally $g'(0)=a/2$, while
$g(z_0)=1$ implies $az_0=0$.
\end{proof}

\begin{remark}[Sharpness]\label{rem:C-algebraic-lucas-gauge}
For $a\ne0$, the cubic gauge $(1-az)^{-1/3}$ has Lucas-prime
density exactly $1/2$. Its degree forces $p\equiv1\pmod3$;
at every sufficiently large prime in that progression its quotient
is $T_p(z)(1-az)^{(p-1)/3}$, of degree at most $5(p-1)/6$.
The theorem therefore has a sharp density threshold. Its extra
normalization selects the standard section among algebraic scalar
gauges; it supplies no complex-to-Frobenius comparison.
\end{remark}

\subsection{Geometric hypotheses in the transport theorem}
\label{sec:countermodel-guardrails}
At a regular algebraic point, the first-jet transformation is
\begin{equation}\label{eq:countermodel-transport}
 \binom{H}{\delta H}
 =g\begin{pmatrix}1&0\\\delta\log g&\delta\log\phi\end{pmatrix}
       \binom{F\circ\phi}{(\thetaop F)\circ\phi},
 \qquad H=gF(\phi).
\end{equation}
Its determinant is $g^2t\phi'/\phi$. The following examples isolate
what fails at a zero of this determinant, at a singular fiber, or
under a change of coefficient field. They concern transformed
families, not the original standard branches.

\begin{proposition}[A critical non-CM pullback]
\label{prop:critical-countermodels}
For an integer $m>3$ with $m\nmid64$, set
\[
 H_m(t)=F_{1/2}((2t-t^2)/m)=\sum_{k\ge0}h_kt^k.
\]
The series and its derivative converge absolutely at $t=1$.
The fiber there is smooth and non-CM, yet
\[
 \delta H_m(1)=0,\qquad
 \sum_{k<p}kh_k\equiv0\pmod{p^3}\quad(p>2,\ p\nmid m).
\]
\end{proposition}
\begin{proof}
Choose $r>1$ with $2r+r^2<m$. The pullback lies in the disk of
convergence on $|t|\le r$, and its derivative at $1$ is zero.
The invariant $64(4m-1)^3/m$ is rational and nonintegral; a CM
invariant is an algebraic integer, so the fiber is non-CM.
For the congruence truncate first in the source coordinate:
\[
 V_p(t)=\sum_{n<p}\frac{\binom{2n}{n}^3}{64^n}
                       ((2t-t^2)/m)^n.
\]
For $n\le(p-1)/2$ each term has degree below $p$; for larger
$n<p$ its coefficient is divisible by $p^3$. Therefore
$V_p-[H_m]_{<p}\in p^3\Z_{(p)}[t]$. As $V'_p(1)=0$ exactly,
differentiation proves the assertion.
\end{proof}
The primitive pair is $(0,1)$ and the multiplier is zero. Here the
proposed de Rham direction itself vanishes; it is not a complement
to the Hodge line. Multiplying by $1+t$ also gives the complex zero
identity with primitive weights $(-1,2)$.

A vanishing gauge destroys uniqueness more completely:
$(1-t)^2F_{1/2}(t/5)$ and its derivative both vanish at $1$,
although its smooth fiber has non-CM invariant $438976/5$.
Every primitive pair then gives zero. A singular fiber creates a
different failure. For $m>2$, the noncritical pullback
$K_m(t)=F_{1/2}(t(1-t)/m)$ has
\[
 K_m(1)=1,\qquad\delta K_m(1)=-1/(8m),\qquad
 (1+8m\delta)K_m(1)=0;
\]
its image is the nodal parameter zero.

\begin{proposition}[Rational coefficients need not preserve rational slopes]
\label{prop:gauge-rational-slope-countermodel}
The normalized rational-coefficient germ
\[
 L(t)=\frac{1+\sqrt{1+t/3}}2F_{1/2}(t/4)
\]
has an absolutely convergent first jet at $t=1$ and a CM fiber there,
but no nonzero rational weight pair gives an inverse-$\pi$ identity.
\end{proposition}
\begin{proof}
Both factors are analytic on a disk of radius greater than one.
At $z=1/4$, the standard identity is
$(1+6\thetaop)F_{1/2}=4/\pi$ and $j=54000$.
Projective uniqueness makes $1/6$ the only algebraic slope.
Since $\delta\log g(1)=(2-\sqrt3)/4$, jet transport gives the unique
slope
\[
                 A/B=-1/3+\sqrt3/4\notin\Q.
\]
For $B=1$ the multiplier is $2g(1)/3\ne0$, so this is an actual
algebraic identity line with no rational point.
\end{proof}

At a singular source value even the chain rule must be interpreted
through analytic germs. Clausen gives
\[
 M(t)=F_{1/2}(t(2-t))={}_2F_1(1/2,1/2;1;t/2)^2,
\]
whose first jet converges at $1$. The Legendre relation at parameter
$1/2$ gives $\delta M(1)=2/\pi$; the fiber has $j=1728$.
Thus the vanishing pullback derivative does not force a zero value
when the source derivative is singular. This also shows how a
pullback can improve ordinary convergence.

\subsection{Boundary terms and finite-jet obstructions}
Raw truncation does not commute with the Euler gauge. Direct
coefficient summation yields
\[
 \left[\frac{F_s}{1-t}\right]_{<p}
       =\frac{T_p(t)-t^pT_p(1)}{1-t}\pmod p.
\]
The boundary term explains the following failure of a raw arithmetic
criterion for a genuine complex identity.

\begin{proposition}[A CM Euler identity with zero density $1/2$]
\label{prop:euler-raw-density-obstruction}
For $H=F_{1/2}/(1-t)$, the identity
$(-1+6\delta)H(1/4)=16/(3\pi)$ has raw weighted-truncation zeros
of natural prime density exactly $1/2$.
\end{proposition}
\begin{proof}
Differentiating the boundary formula in characteristic $p$ gives
\[
 (-1+6\delta)[H]_{<p}(1/4)
 \equiv\tfrac43(1+6\thetaop)T_p(1/4)-\tfrac13T_p(1).
\]
The first term vanishes outside finitely many primes by the standard
CM criterion. Finite Clausen identifies the second with
$U_p(1/2)^2=P_{(p-1)/2}(0)^2$. For a Legendre polynomial,
$P_m(0)=0$ when $m$ is odd and
$P_m(0)=(-1)^{m/2}2^{-m}\binom m{m/2}\ne0\pmod p$
when $m$ is even. Thus the zero primes are, up to finitely many,
exactly $p\equiv3\pmod4$.
\end{proof}

\label{subsec:C-gauge-obstruction}
\begin{proposition}[Finite jets do not determine prime-zero density]
\label{prop:C-finite-jet-density-zero}\label{cor:C-two-jet-density-zero}
Let $M\ge2$, $N\ge1$, $F=F_{1/2}$, and
\[
 g_{M,N}(z)=1+\frac{z^N(4z-1)^M}{3-4z},\qquad H=g_{M,N}F.
\]
Then $H(192t)\in\Z[[t]]$, and the minimal differential module of
$H$ is isomorphic over $\Q(z)$ to the irreducible rank-three module
of $F$. The difference $H-F$ vanishes to order $N$ at zero and
order $M$ at $1/4$. Hence the ordinarily absolutely convergent
identity $(1+6\thetaop)H(1/4)=4/\pi$ holds. Nevertheless the primes
at which $(1+6\thetaop)[H]_{<p}(1/4)=0$ have natural density zero.
\end{proposition}
\begin{proof}
Put $P=3-4z+z^N(4z-1)^M$. For $p\ge2(M+N)+1$,
$\deg(PT_p)<p$, so finite geometric summation gives
\begin{equation}\label{eq:C-gauge-finite}
 [H]_{<p}(z)=
 \frac{P(z)T_p(z)-(4z/3)^pP(3/4)T_p(3/4)}{3-4z}
 \quad\text{in }\mathbb F_p[z].
\end{equation}
The quotient is polynomial because its numerator vanishes at $3/4$;
its germ agrees with $H$ through degree $p-1$. At $1/4$ the first
jets of $H$ and $F$ agree. Frobenius kills the derivative of
$(4z/3)^p$, while
$(1+6\thetaop)(3-4z)^{-1}|_{z=1/4}=2$. Thus
\begin{equation}\label{eq:C-gauge-weighted}
 (1+6\thetaop)[H]_{<p}(1/4)
 =(1+6\thetaop)T_p(1/4)
       -2^{M+1-2N}3^{N-1}T_p(3/4).
\end{equation}
For $N=1$ the coefficient is $2^{M-1}$. The first term is cofinally
zero by CM. Finite Clausen identifies $T_p(3/4)$ with the square
of the Hasse polynomial of the Legendre curve at parameter $1/4$.
Its invariant is $35152/9$, so this curve is non-CM. Serre's
open-image theorem and Chebotarev, in the form of
Theorem~\ref{thm:B-qcurve-density}, give supersingular-prime density
zero. This proves the density assertion.

The remaining properties express elementary functoriality. Rational
multiplication by $g_{M,N}\ne0$ is an invertible change of differential
module; the numerator gives the two vanishing orders; and all factors
are analytic on $|z|<3/4$. Finally $F(192t)$ has integral coefficients,
and the denominator of the added gauge term becomes $3(1-256t)$,
whose factor $3$ divides $192^N$. Hence the stated integral rescaling
holds.
\end{proof}
The construction preserves any prescribed finite jets at both the
cusp and the identity point. It changes the generating function,
while retaining its CM fiber; it is not a counterexample to standard
CM forcing.

\begin{proposition}[Finite jets with Lucas congruences on a progression]
\label{prop:C-algebraic-jet-lucas}
For $M\ge2$, $N\ge1$, put $D=M+N$ and choose an integer $m\ge2D$.
The branch
\[
 g(z)=\bigl(1+2^{-M-1}z^N(4z-1)^M\bigr)^{-1/m},\qquad g(0)=1,
\]
is nonrational. For $H=gF_{1/2}$, the difference $H-F_{1/2}$
vanishes to orders at least $N$ at zero and $M$ at $1/4$;
$(1+6\thetaop)H(1/4)=4/\pi$ converges absolutely. Moreover $H$
is Lucas at every sufficiently large prime $p\equiv1\pmod m$.
\end{proposition}
\begin{proof}
Writing the polynomial inside parentheses as $P$, one has
$|P-1|\le2^{-1}4^{-N}<1$ on $|z|\le1/4$.
The normalized root therefore extends beyond that disk and takes
value one at $1/4$. Its vanishing orders prove the complex assertions.
It cannot be rational, since its order at infinity would imply
$m\mid D$, whereas $0<D<m$. At a good prime $p\equiv1\pmod m$,
\[
 \mathcal Q_p(H)=T_pP^{(p-1)/m},\qquad
 \deg\mathcal Q_p(H)=\tfrac{p-1}{2}+D\tfrac{p-1}{m}\le p-1.
\]
The polynomial quotient criterion proves Lucas.
\end{proof}
The progression restriction is essential by the half-density theorem.
Together these results identify arithmetic normalization as additional
data that finite complex jets cannot recover. They leave open the
comparison of the standard complex identity with exact Frobenius
preservation.
